\chapter{A matrix action from ordered diagonal operations}
\label{chap:cmc-matrix-action}

The cyclic model of Chapter~\ref{chap:cyclic-mixed-diagonal} gives a finite
action which retains its higher ordered products. We now construct that
action on actual matrix coordinates and verify its master equations.

\section{The finite matrix action}
\label{sec:cmc-hamiltonian}

The cyclic pairing packages the operations into an action.
To retain their order, first use the free graded algebra on
even letters \(x,u\) and odd letters \(v,t\), completed in word
length. No relation \(v^2=0\) or \(t^2=0\) is imposed in this
free algebra. On the fibre \(y=z=0\), introduce an even
coupling \(\lambda\). Multiplying the \(n\)-ary operations by
\(\lambda^{n-2}\) for \(n\ge3\) preserves their identities:
every composite in an identity has the same total power of
\(\lambda\). The binary products are unchanged.

Define the odd derivation \(d\) on \(x,u,v\) and its even
curvature \(h\) by
\begin{align}
 dx={}&\lambda(-xuv+xvu+uxv+uvx+vxu-vux)
          -\lambda^2(uv^3+v^3u),\notag\\
 du={}&\lambda(-x^2v-xvx-vx^2-u^2v+uvu-vu^2)
          +\lambda^2(xv^3+v^3x),\notag\\
 dv={}&ux-xu,\label{eq:cmc-curved-d}\\
 h={}&\lambda(x^3+xu^2-uxu+u^2x)\notag\\
 &-\lambda^2(x^2v^2+vx^2v+v^2x^2+
                  u^2v^2+vu^2v+v^2u^2)+\lambda^4v^6.
 \label{eq:cmc-curvature}
\end{align}
These satisfy \(d^2a=ha-ah\) and \(dh=0\).
One can derive them by dualizing the cyclic table. A tensor
evaluation on a word with \(m\) odd dual letters has sign
\((-1)^{m(m-1)/2}\), and the differential of a homogeneous
functional \(f\) is \(-(-1)^{|f|}f b\). Including the dual of
the unit gives the odd derivation
\[
 Qt=-t^2+h,\qquad
 Qa=-[t,a]_{\mathrm s}+da\quad(a=x,u,v).
\]
The table identity gives \(Q^2=0\); substituting these unit
terms gives precisely \(d^2=[h,-]\) and \(dh=0\).
Equivalently, expand the three generator identities in free
words using~\eqref{eq:cmc-curved-d}. Equality on the generators
is sufficient, since both sides are even derivations.

Let now \(X,U\) be even and \(V,T\) odd \(N\)-by-\(N\)
matrices, with entries in a graded-commutative polynomial
algebra. Odd entries anticommute; products of odd matrices
need not square to zero. Matrix trace has the graded cyclic
rule. Define the odd function
\begin{align}
 \mathcal H_N=\operatorname{Tr}\bigl(
 &-VT^2+T(UX-XU)\notag\\
 &+\lambda V(X^3+XU^2-UXU+U^2X)\notag\\
 &-\lambda^2V^3(X^2+U^2)+\lambda^4V^7/7\bigr).
 \label{eq:cmc-matrix-action}
\end{align}
Define an even Poisson bracket by the generator values
\[
 \{X_{ij},U_{kl}\}=-\delta_{il}\delta_{jk},\qquad
 \{T_{ij},V_{kl}\}=\delta_{il}\delta_{jk}.
\]
The reverse values follow from
\(\{A,B\}=-(-1)^{|A||B|}\{B,A\}\); every other generator
pair has bracket zero. Extend by linearity and
\(\{A,BC\}=\{A,B\}C+(-1)^{|A||B|}B\{A,C\}\).
The Jacobi identity holds on generators, whose brackets are
constant, and hence on all polynomials by this product rule.
Its nondegenerate coefficient matrix is the inverse of an even
symplectic form. The Hamiltonian derivation
\(Q=\{\mathcal H_N,-\}\) has equations
\[
 QT=\partial_V\mathcal H_N,\quad
 QV=\partial_T\mathcal H_N,\quad
 QX=\partial_U\mathcal H_N,\quad
 QU=-\partial_X\mathcal H_N.
\]
The derivatives are cyclic derivatives, including the transpose
of matrix indices. Moving a prefix of an odd cyclic word to
its end has sign \(+1\), since its parity \(p\) gives
\(p(1-p)=0\). Differentiation of~\eqref{eq:cmc-matrix-action}
therefore gives exactly the matrix evaluations of the preceding
four formulas for \(Q\). For instance differentiating
\(-V^3X^2\) with respect to \(V\) gives
\(-V^2X^2-VX^2V-X^2V^2\).

Thus \(Q^2=0\) at every finite rank. The Jacobi identity gives
\(Q^2(f)=\tfrac12\{\{\mathcal H_N,\mathcal H_N\},f\}\).
The Hamiltonian of its
square is \(\{\mathcal H_N,\mathcal H_N\}/2\); nondegeneracy
of the symplectic form makes this function constant. Its
polynomial order is positive, so the constant is zero.
This is the classical master equation in this finite
symplectic space.

There is also a finite cotangent BV action. List the matrix
entries as coordinates \(\phi^i\), of parities \(\epsilon_i\),
and introduce dual coordinates \(\phi_i^*\) of parities
\(\epsilon_i+1\). Use the canonical odd bracket
\(\{\phi^i,\phi_j^*\}_{\mathrm{BV}}=\delta^i_j\), with all
base--base and dual--dual brackets zero. Extend it by
\[
 \{F,G\}_{\mathrm{BV}}
 =-(-1)^{(|F|+1)(|G|+1)}\{G,F\}_{\mathrm{BV}},
\]
and by
\[
 \{F,GH\}_{\mathrm{BV}}
 =\{F,G\}_{\mathrm{BV}}H
   +(-1)^{(|F|+1)|G|}G\{F,H\}_{\mathrm{BV}}.
\] With left derivatives put
\[
 S=\sum_i(-1)^{\epsilon_i}\phi_i^*Q\phi^i,\qquad
 \Delta=\sum_i(-1)^{\epsilon_i}
             \partial_{\phi^i}\partial_{\phi_i^*}.
\]
The action \(S\) is even and
\(\{S,\phi^i\}_{\mathrm{BV}}=Q\phi^i\).
On the dual variables the Hamiltonian derivation acts by the
dual derivative of \(Q\). Differentiating \(Q^2=0\) shows that
its square also vanishes there. The bracket identity therefore makes \(\{S,S\}_{\mathrm{BV}}\)
commute with every coordinate. Its partial derivatives vanish,
so in characteristic zero it is a coefficient constant. Every
term still contains one dual coordinate, and therefore vanishes
on the zero section. Thus \(\{S,S\}_{\mathrm{BV}}=0\).
Moreover
\[
 \Delta S=\sum_i\partial_{\phi^i}(Q\phi^i)=0.
\]
For each even pair \(X,U\), the two mixed derivatives cancel.
For each odd pair \(T,V\), they cancel because the two left
derivatives anticommute. The matrix transpose in each pairing
matches the corresponding summation indices. This is the
zero-divergence identity in these explicit conventions.
It follows that
\[
 \frac12\{S,S\}_{\mathrm{BV}}+\hbar\Delta S=0
\]
for formal \(\hbar\).
This finite polynomial statement supplies an action and all
of its identities. It does not introduce a spatial propagator,
boundary conditions, or a renormalized field theory.

At rank one, the odd scalar squares vanish and even
coordinates commute, so the action reduces to
\(\lambda v(x^3+xu^2)\).
The longer ordered terms are exactly the terms this reduction
forgets. Recovering the matrix theory from that scalar
potential alone would lose its noncommutative operations.

\section{Retaining the base point}

The full parameter-dependent model also has a matrix action.
Let \(\mathcal H_N^0\) denote
\eqref{eq:cmc-matrix-action} with \(\lambda=1\). Define
\[
 \mathcal H_N^{y,z}
 =\mathcal H_N^0(X+yI,U+zI,V,T)
       -(y^3+yz^2)\operatorname{Tr}V.
\]
Translation in the even coordinates preserves the bracket.
Furthermore
\[
 \{\mathcal H_N^0,\operatorname{Tr}V\}
 =\operatorname{Tr}(UX-XU-TV-VT)=0,
\]
by graded trace cyclicity, and
\(\{\operatorname{Tr}V,\operatorname{Tr}V\}=0\).
The subtraction therefore preserves the master equation.

This is the action of the full relative cyclic table, not only
of its generator restriction. To verify the identification,
apply its graded dual rule before evaluating matrices.
The coefficient of each ordered word in the four resulting
derivations agrees with translation \(x\mapsto x+y\),
\(u\mapsto u+z\) in the fibre formulas. In \(Qt\) subtract
the scalar \(y^3+yz^2\). For example its linear coefficients
are \((3y^2+z^2)x+2yz\,u\), dual to the unary differential;
its coefficient of \(v^2\) is \(-3(y^2+z^2)\), dual to
\(\widehat b_2(c,c)\) with its tensor sign. The ordered
coefficients in the remaining rows give the remaining
terms by the same rule.

At rank one the relative Hamiltonian is
\[
 v\bigl(W(x+y,u+z)-W(y,z)\bigr).
\]
The full matrix formula retains the ordered terms that this
commutative restriction removes.

Likewise the cyclic form constructed here does not by itself
identify a specified geometric negative-cyclic class: that
identification requires a comparison of its full cochain
representative, beyond the leading bilinear pairing.
